\section{Three weight classes and the remaining degree-18 datum}
\label{sec:cubic-three-weights}

The full module first stabilizes at $E_2$, as AC-K20 proves.
Its positive-weight specializations need less information.
The same certified first-layer bound used in AC-K18
forces the following sharper result.

\begin{theorem}[AC-K21: exponents four and five with period three]
\label{AC-K21}
For the fixed cubic character over $K=\mathbb Q(\omega)$,
\[
          4\leq m_e\leq5,
       \qquad m_{e+3}=m_e\quad(e\geq1).              \tag{KJ.1}
\]
There are only three possible sequences:
\[
 \begin{array}{c|ccc}
 &e\equiv1\pmod3&e\equiv2\pmod3&e\equiv0\pmod3\\\hline
 \text{first possibility}&4&4&4\\
 \text{second possibility}&5&4&4\\
 \text{third possibility}&4&5&4.
 \end{array}                                         \tag{KJ.2}
\]
Equivalently, either every $m_e=4$, or there is one
$r\in\{1,2\}$ such that $m_e=4+\mathbf1_{e\equiv r\ (3)}$.
For $N=t-1\geq1$, the exact aggregate in these alternatives is
\[
 \sum_{e=1}^Nm_e=
 \begin{cases}
 4N,&\text{in the first possibility},\\
 4N+\left\lfloor (N+3-r)/3\right\rfloor,
                  &\text{in the possibility indexed by }r.
 \end{cases}                                         \tag{KJ.3}
\]
In particular the fixed coefficient-prime-three budget is
at most $(4N+\lceil N/3\rceil)\log3$.
All the wild kernels $W_e(E)$ are killed by nine.

The choice among (KJ.2) is determined by a specified
conjugation component of
\[
 C_1=\widetilde{\operatorname{Cl}}_3
        (\mathbb Q(\mu_9,\sqrt[3]{269})),             \tag{KJ.4}
\]
with its normalized cyclotomic action.  The criterion below
is finite and uses the degree-18 field; this theorem does
not compute which alternative occurs.
\end{theorem}
\begin{proof}
\proofstep{1. Specializing the cyclic plus component.}
Keep $M=X'$, $R=\mathbb Z_3[[B]]$, $B=\gamma-1$,
and $M_+=Rx$, $U=M_-=R\delta x$ from AC-K18.
Here the plus sign is the conjugation sign of $x$;
it need not be the numerical sign $+1$.
The already proved facts are
\[
 \begin{gathered}
 \delta M_+=U,\quad \delta U=-3M_+,\quad
 M_+/BM_+\simeq\mathbb Z/9,\quad U/BU\simeq\mathbb Z/3,\\
 M_+/3M_+\simeq\mathbb F_3[B]/B^d,
       \qquad d\in\{6,7\},\qquad \#U\leq9,\qquad27M=0.
 \end{gathered}                                      \tag{KJ.5}
\]
Write $M_+=R/I$ using the generator $x$.
Reduction modulo three and evaluation at $B=0$ give
\[
       I\bmod3=(B^d),\qquad
       \{f(0):f\in I\}=9\mathbb Z_3,
       \qquad27\in I.                               \tag{KJ.6}
\]
Consequently, for $f(B)=\sum_{k\geq0}f_kB^k\in I$,
$9\mid f_0$ and $3\mid f_k$ for $k<d$.
If $b\in3\mathbb Z_3$, the convergent evaluation therefore
satisfies
\[
                         9\mid f(b).                \tag{KJ.7}
\]
The constant and linear terms are divisible by nine;
the terms of degree at least two are divisible by nine
already from the powers of $b$.
Formal division by $B-b$, valid for $b\in3\mathbb Z_3$,
identifies the specialized plus component with
\[
 M_+/(B-b)M_+
       \simeq\mathbb Z_3/\operatorname{ev}_b(I).
                                                               \tag{KJ.8}
\]
The image here is an ideal, since evaluation is onto
$\mathbb Z_3$; it contains $27\mathbb Z_3$ and lies
inside $9\mathbb Z_3$ by (KJ.6)--(KJ.7).
Thus its quotient has length at least two.

For a positive weight $e$, put $b_e=4^{-e}-1$.
The relation $4^e\gamma-1=0$ is equivalent to $B=b_e$,
and $b_e\in3\mathbb Z_3$.
The image of $x$ generates $W_e=M/(4^e\gamma-1)M$
over $\mathcal O$, because $B$ acts as a scalar there.
By AC-K12 it is a cyclic module of length $a_e=m_e-1$.
Using the original conjugation on $M$, the lengths of
its two components are
\[
       \left\lceil a_e/2\right\rceil
       \quad\hbox{and}\quad
       \left\lfloor a_e/2\right\rfloor.             \tag{KJ.9}
\]
This follows directly from
$\mathcal O=\mathbb Z_3\oplus\delta\mathbb Z_3$
and $\delta^2=-3$.
These are the labels inherited from $M$, irrespective
of the extra sign $(-1)^e$ in the action of actual
conjugation on the Tate twist.
The lower bound two in (KJ.8) implies $a_e\geq3$.

\proofstep{2. Every specialized plus component is killed by nine.}
If $U$ is elementary abelian, then $3U=0$ and
$3M_+=-\delta U$ give $9M_+=0$.
Together with (KJ.8), this makes every specialized
plus component have length exactly two.

The remaining possibility is $U\simeq\mathbb Z/9$.
By (KJ.5), $B$ acts on $U$ as multiplication by
$b_U=3c$ modulo nine, where $c\in\mathbb F_3^\times$.
Set $F(B)=(1+B)^3-1$.
This endomorphism kills $U$.
Logarithmic descent and the certified first-layer bound
AC-K15 give
\[
 C_1=M/F(B)M=C_+\oplus U,
 \qquad \operatorname{length}_{\mathbb Z_3}C_+\leq4,
 \qquad C_+=M_+/F(B)M_+.                             \tag{KJ.10}
\]
Modulo three, $F(B)=B^3$; hence $C_+/3C_+$ has
dimension three, since $d\geq6$.
Also $C_+/BC_+=M_+/BM_+\simeq\mathbb Z/9$.
Thus $C_+$ is not elementary abelian.
Its length bound and its three generators modulo three
force
\[
                   C_+\simeq\mathbb Z/9\oplus
                                     (\mathbb Z/3)^2.
                                                               \tag{KJ.11}
\]
In particular there exists $h(B)\in R$ such that
\[
                       9x=F(B)h(B)x.                \tag{KJ.12}
\]
Reduce this identity modulo three in
$M_+/3M_+=\mathbb F_3[B]/B^d$.
It gives $B^3h(B)=0$ in that quotient, and therefore
\[
                       h(B)\in(3,B^{d-3}).          \tag{KJ.13}
\]
For every $b\in3\mathbb Z_3$, this implies
$3\mid h(b)$; also
$F(b)=3b+3b^2+b^3$ is divisible by nine.
After specialization in (KJ.12), its right side is
therefore a multiple of $27x$, which is zero by
(KJ.5).  Thus $9x=0$ in every specialized plus
component.  Equation (KJ.8) again makes its length
exactly two.

\proofstep{3. The remaining component and the three possibilities.}
If $U$ is elementary abelian, then $4^e=1$ on $U$,
and its quotient by $4^e\gamma-1$ is $U/BU$ of
length one.  The plus length is two, so $a_e=3$
and $m_e=4$ at every weight.
If $U=\mathbb Z/9$ with $B=b_U=3c$ as above, then
\[
                  4^e(1+3c)-1\equiv3(e+c)\pmod9.
                                                               \tag{KJ.14}
\]
Its quotient has length two for $e\equiv-c\pmod3$
and length one elsewhere.
The plus length is always two, so these alternatives
give respectively $a_e=4$ and $a_e=3$, proving
(KJ.1)--(KJ.2).  In particular $a_e\leq4$;
since $(\lambda^4)=(9)$, nine kills every $W_e$.
Counting the integers in the exceptional residue class
gives (KJ.3).
\end{proof}

\paragraph{The exact degree-18 criterion.}
The proof of AC-K20 shows that $F(B)U=0$ also in the
elementary cases.  Consequently the component of $C_1$
with sign opposite to that of $x$ is exactly $U$, with
its cyclotomic action.  The sign of $x$ can be identified
intrinsically: it is the unique conjugation sign whose
component of
$C_1/BC_1=M/BM\simeq\mathbb Z/9\oplus\mathbb Z/3$
has order nine.  If this opposite component is elementary,
all $m_e=4$.  If it is cyclic of order nine, compute
$\gamma=1+3c$ on it using the field automorphism
$\zeta_9\mapsto\zeta_9^4$, $\sqrt[3]{269}$ fixed.
Then $r\equiv-c\pmod3$ is the unique class with $m_e=5$.
Thus a certificate for just this component and its action
would choose among (KJ.2); a complete certified
presentation of $C_1$ with conjugation and cyclotomic
actions is sufficient.  The first-layer upper certificate
AC-K15 does not by itself supply these exact actions or
select an alternative.

The degree-54 reconstruction of the whole module in AC-K20
remains valid, as does its nine-quotient recipe.  For the
exponents alone the present theorem reduces that recipe
to three specializations and, more economically, to the
degree-18 component just specified.  The finite-cover and
even K-theory interpretation following AC-K20 applies
with these sharper integral lengths.  No new external
arithmetic theorem, numerical certification, or archive
locator enters this finite-module deduction.
